% !TEX root = ../thesis.tex
\chapter{Konstrukce ASR systému pro uživatele po totální laryngektomii hovořící pomocí elektrolarynxu}
\label{chap:construction}


% v~předchozích sekcích bylo zmíněno, že velmi významný problém představuje psychologický faktor ztráty hlasu. Ten sebou nese i na první pohled néúplně zřejmou komplikaci. A tou je obtížnost získávání řečníků po TL, kteří jsou ochotní spolupracovat na výzkumu.

% Jedním z možných řešení je, že se výzkumníci sami naučí používat jednu z výše  \todo{TBD}{[c4l1]})zmíněných metod produkce řeči. Pro účely našeho výzkumu a rámec této práce by se jednalo o elektrolarynx. I když se tato možnost jeví poměrně jednoduše, tak zdání klame. Přestože je zdravý řečník schopen obstojně mluvit (pomocí EL), za relativně krátkou dobu tak k tomu, aby produkovaná řeč svými parametry odpovídala zkušenému řečníkovi po TL, je potřeba relativně dlouhá doba.ě

% v~našem případě se podařilo získat pouze jednoho zkušeného řečníka\footnote{Jedná se o ženu v~důchodovém věku, která ale stále působí na akademické půdě a jednou za čas i přednáší.}, který podstoupil TL před více než 15 lety a EL používá aktivně každý den. Je to jeho jediná možnost jak produkovat slyšitelnou řeč.

% Díky spolupráci s tímto řečníkem jsme byli schopni získat cca 15 hodin promluv (více o získaných datech v~\ref{chap:experiments:analysis} a \ref{chap:experiments:normalization}), které byly použity pro všechny dosavadní experimenty. Z pohledu standardních obecných systémů rozpoznávání řeči se to může jevit jako velmi malé množství, ale jak ukáží následné experimenty, tak to není velký problém. Jelikož už před prvními experimenty se jevilo jako prozaičtější vytvářet individuální modely pro každého řečníka a výsledky experimentů toto jen podpořily.

% Při rozhovorech s řečníkem se potvrdil psychický aspekt ztráty hlasu na člověka po TL. Konkrétní osoba ještě dlouhá léta po operaci nebyla schopna telefonovat, natož mluvit na veřejnosti. Kvalita života se tímto velmi snížila a trvalo prý velmi dlouho dobu, než se daná osoba odvážila i jen odpovědět na nečekaný telefonní hovor. To pro nás představovalo nezanedbatelnou porci motivace v~práci.

% v~následujícím textu si přiblížime a zanalyzujeme pořízená data (\ref{chap:experiments:analysis}) a na základě těchto prvotních experimentů nadefinujeme cíle a určíme další postup. v~části (\ref{chap:experiments:normalization}) si popíšeme normalizaci dat a porovnáme schopnosti člověka a stroje. v~sekci \ref{chap:realisation:augmentation} rozebereme prvotní ověřovací experimenty s upravenými daty a v~sekci \ref{chap:experiments:durationmodels} představíme výsledky modelů zohledňujících i délku fónému.

% Při konzultacích s pacienty vyšlo najevo, že psychická stránka je velmi důležitá a někdy to může věst až k paralýze (v obrazném slova smyslu) kdy ti lidé nebyli schopni třeba telefonovat. Takže systém, který by sice neřešil kompletní problematiku náhrady/rehabilitace řeči by mohl být užitečný.

% !TEX root = ../thesis.tex
\section{Vytvoření řečového korpusu EL promluv}
\label{chap:construction:corpus}

Na začátku byla idea o pomoci skupině lidí mající problémy s přirozenou řečí. Vůbec prvním předpokladem, na cestě k úspěšnému dosažení tohoto cíle, jsou data. Jelikož se jedná o velmi spefická data, tak je potřeba zajistit co možná největší množství kvalitních\footnote{Kvalitou je myšlena věrnost dat dané doméně, dále se mluví o přesnosti ve smyslu bezchybnosti přepisů.} a přesných dat.

V části \ref{chap:cause:desease} bylo zmíněno, že ročně se objevý více než 100 nových případů trvalé ztrázy hlasu ročně. Zaroveň bylo řečeno \cite{Skvrnakova2010}, že více rizikovými osobami jsou starší lidé, kteří intenzivně kouří a konzumují alkohol. Přesto je patrný trend snižujícího se věku pacientů a s tím souvisejícím nárůstem případů ztráty hlasu. Přičteme-li již zmíněný psychologický aspekt jeho ztráty, tak je zřejmé, jak komplikované je získat ke spolupráci i jen jednoho řečníka ochotného podstoupit naročné\footnote{I pro zdravého člověka je někdy někalikahodinové nahrávání vysilující. Pro jedince po TL to je z mnoha důvodů ještě řádově náročnější.} nahrávání.

Při libovolné práci s pacienty po TL, dřív nebo později dojde k určité formě spolupráce s oddělením ORL, které má nastarosti péči o tyto pacienty. V našem připadě nejprve s ORL klinikou při Fakultní nemocnici v Plzni a poté i s ORL klinikou Fakultní nemocnice v Motole. S jejich pomocí jsme získali ke spolupráci jednoho řečníka. Konkrétně se jedná o dámu v duchodovém věku, která podstoupila TL před více než 15 lety. Po překonání ostychu\footnote{Podle jejích vlastních slov nebyla schopna několik let po operaci ani zvednout nečekaný telefonní hovor, natož mluvit na veřejnosti.} se byla schopna naplno vrátit do běžného života a dokonce v určité formě opět přednášet o stomatologii na Lekařské fakultě v Plzni Univerzity Karlovy.

S její pomocí jsem, v 1. etapě nahrávání, byli schopni pořídit přes 10 hodin promluv, viz tab. \ref{tab:construction:recording}. Nahrávání probíhala v relativně spartánských podmínkách za plného běžného provozu katedry. Přesto získaná data neosahují žádný nežadoucí ruch, kromě toho produkovaného samotným EL.

Nahravací aparatůra sestávala z miniaturního profesionálního mikrofonu (DPA d:screet 4061-FM), zesilovačem (DPA MMA6000), externí zvukovou kartou a běžného notebooku. Mikrofon byl pomocí bezpolštářkové náplasti přilepen poblíž pravého koutku úst, abychom zaznamenaná řeč měla co možná nejvyšší kvalitu.

Celé nahrávání bylo v 1. etapě rozděleno do 14 samostatných sezení a probíhalo od prosince roku 2010 do května roku 2011. Každé sezení trvalo přibližně dvě hodiny během kterých se podařilo získat necelou hodinu akustických dat. Samotné nahrávání se sestávalo z 10 - 20 minutového úseku pořizování nahrávky a přibližně 10 minut dlouhého odpočinku. Ten byl nezbytný hlavně z důvodu únavy řečníka.

Ještě před samotným nahrávánám byly pečlivě vybrány a vytvořeny 2 sady vět:

\begin{enumerate}
  \item sada obsahující všechny možné české fonémy - \textit{40 vět}.
  \item sada obsahující věty s reálnou četností fonémů - \textit{5000 vět} \cite{Radova2000}.
\end{enumerate}

\noindent Pořízené nahrávky vždy odpovídají 10 - 20 minutovému úseku nepřerušovaného nahrávání a soubory tak vždy obsahují několik vět. Ty jsou od sebe odděleny minimálně 5 sekundovým úsekem ticha. Nahrávky dále mouhou obsahovat opakování chybně vyslovené věty, přeřeknutí, kýchnutí a další neřečové události. Z tohoto důvodu bylo nezbytné pořízené nahrávky anotovat, přestože byly pořízené na základě připravené sady vět.

Ještě před samotným anotováním byly nahrávky, podle úseků s tichem, rozsekány na menší části. V tomto případě úplně dobře nefungovali\footnote{Problémem byl zvuk EL, se kterým nebylo při návrhu VAD počítáno.} standardně používané sofistikovanější metody pro voice activity detection (angl. zkratka VAD), a proto bylo využito principu energie. Pro každou nahrávku obsahující více vět se pomocí vzorce

\begin{equation}
  \label{eq:construction:energy}
  E_{RMS}(n) = \sqrt{\frac{1}{N} \sum_{n=1}^{N} \left| x(n) \right|^2},
\end{equation}

\noindent kde $N$ představuje počet vzorků v nahrávce a $x(n)$ představuje pravoúhlé okénko vzorku $n$. Pro tento případ se ukázalo jako vhodnější volit root-mean-square energy ($E_{RMS}$) a empericky se ukázalo, že vhodná délka okénka je v rozmezí $10 - 100$ ms. Na obr. \ref{fig:construction:el_speech} je zobrazena podoba audio signálu a spektrogram promluvy \textit{\uv{Akcie Komerční banky}}. Zároveň je zde vypočtené hodnoty energie a celková průměrná energie. Tyto hodntoty slouží pro určení míst kde začíná a končí věta. Na začátku a konci každé věty je dobré mít minimálně $0.5$ s ticha. Tím pádem, pokud energie nějakého úseku $x$ je $E_{RMS}(x) < avg(E_{RMS})$ a zároveň délka tohoto úseku $dur(x) \geq 1\ [s]$, tak nahrávku můžeme v tomto úseku rozdělit.

\begin{figure}[hbpt]
  \centering
  \includegraphics[width=0.9\textwidth]{./ch5-construction/img/energy_spec_el.png}
  \caption{Průběh a spektrogram promluvy a vyznačenou energií EL promluvy.}
  \label{fig:construction:el_speech}
\end{figure}

Samozřejmě pokud řečník v průbehu věty z libovolného důvodu udělal větší pauzu než $1$ s, tak tato věta byla rozdělena. Jelikož jsou výsledné kratší useky promluv následně anotovány, tak to nepředstavuje problém. Pro budoucí zpracování není podstatné zda promluva je opravdu celá věta, ale to jestli je tento úsek správně přepsán. Fakt, že některé věty jsou rozděleny je důvodem proč v tab. \ref{tab:construction:recording} více souborů než vět.

K anotaci posloužil interní nástroj určený k tomuto účelu a podíleli se na něm celkem 3 anotátoři z řad studentů, kteří si vzájemně kontrolovali své anotace. Ačkoli bylo potřeba anotovat relativně malé množství dat (cca 10 hodin audio záznamu), tak anotace zabrala přibližně 2 měsice. Hlavním důvodem byla relativně dlouhá doba, po kterou se anotátoři adaptovali na specificka EL řeči. Problémem bylo to, že nejprve nebyli vůbec schopni poruzumnět obsahu promluvy a tím pádem jej správně přepsat.

Pokud je pro produkci řeči použit elektrolarynx, tak vedlejším produktem je nezanedbatelný ruch způsobený samotným zařízením \ref{chap:cause:treatment:foniatric}. Přeci jen jeho jedinout funkcí je vybudit vzduch v dutině ústní a tím umožnit produkci slyšitelné řeči. Z tohoto důvodu byly v průběhu anotace ignorovány v podstatě všechny skupiny neřečových událostí, protože vetšinu obsahu nahrávek by bylo nezbytné anotovat jako, že obsahují šum.

Výsledný korpus tedy představuje $5040$ unikátních vět rozdělených do $6385$ souborů (viz tab. \ref{tab:construction:recording}), které v průměru obsahují $7$ slov o průmerné délce $5$ znaků. Tento korpus slouží jako základ pro všechny budoucí experimenty.

\begin{table}[htpb]
  \centering
  \def\arraystretch{1.5}
  \pgfplotstabletypeset[
    col sep=comma,
    string type,
    columns/phase/.style={column name=Nahrávání, column type={|l}},
    columns/length/.style={column name=Délka \textit{[HH:MM:SS]}, column type={|r}},
    columns/sentences/.style={column name=Počet vět, column type={|r}},
    columns/files/.style={column name=Počet souborů, column type={|r|}},
    every head row/.style={after row=\hline, before row=\hline},
    every last row/.style={after row=\hline},
  ]{./ch5-construction/tabs/02-recording1-stats.csv}
  \caption{Informace o korpusu nahrávek z 1. etapy nahravání.}
  \label{tab:construction:recording}
\end{table}

% !TEX root = ../thesis.tex
\section{Analýza akustického signálu a jeho parametrizace}
\label{chap:construction:analysis}

Rozpoznávání řeči se věnuje nemalé usílí již od 50. let 2O. století a v současné době nikoho nepřekvapí téměř bezchybně fungující obecný rozpoznávač v mobilních zařízeních. Pro obecné systémy dokonce existují korpusy s desítkami či stovkami a více hodin promluv, které je možné využít při vytváření těchto systémů.

Tyto korpusy však obsahují ve většině případů pouze \uv{standardní}\footnote{Slovením spojením \uv{standardní řeč} je myšlena řeč neobsahující vyrazné řečové vady, případně jiné formy produkce a často v nepřílíš akusticky náročném prostředí.} řeč. Pokud je snaha vytvořit nebo ověřit funkčnost systému za specifických podmínek (ať už se jedná o rušné prostředí či speciální typy promluv), tak je nezbytné získat potřebná data.

\subsection{Analýza získaných dat}
\label{chap:construction:analysis:data}

Po dokončení anotace obsahuje korpus přes 10 hodin akustických záznamů promluv a více či méně přesných přepisů\footnote{I přes nemalou snahu a několikastupňovou kontrolu, je téměř jisté, že by nebylo obtížné najít přepis, který obsahuje chybu například ve formě překlepu.}. Když jsou k dispozici data je možné se podívat na specifika EL řeči a případně porovnat se zdravým řečníkem.

Pro potřeby porovnání byl použit začátek promluvy \textit{\uv{Akcie Komerční banky...}}. Tuta promluva je součástí standardní množiny vět používaných při vytváření řečových korpusů na KKY při ZČU. Tím pádem je k dispozici v relativně velkém množství příkladů pro zdravé řečníky a také je součástí korpusu EL řeči.

Na obr. \ref{fig:construction:el_speech} a \ref{fig:construction:normal_speech} je zobrazen průběh signálu a spektrogram vybrané promluvy. Už na první pohled je možné zaznamenat určité rozdíly. Prvním takovým je délka promluvy, v případě zdravého řečníka je o celou 1 vteřinu kratší než v případě EL řeči. Tempo řeči je samozřejmě velmi individuální, ale z principu je EL řeč pomalejší. Z průbehu signálu na obr. \ref{fig:construction:el_speech} je patrné, že řečník dělá výraznější pauzy mezi jednotlivými slovy promluvy. To může být způsobené například potřebou naplnit jícen vzduchem. Po TL je dýchání realizováno přes tracheu a pokud nebyl voperován shunt (více v \ref{chap:cause:treatment:tracheo}), tak je trvale oddělen hrtan a hltan. Přesto, pro produkci některých neznělých fonémů je potřeba exhalovat vzduch z dutiny ústní. Zkušený EL řečník to dělá naprosto automaticky, nicméně \uv{polykání} vzduchu zabere nějaký čas. Nevyhnutelným důsledkem je pak velmi častý výskyt samovolného říhání v průběhu promluvy\footnote{Fakt, že je říhání jako neřečová událost běžnou součástí téměř každé promluvy, vedl k ignorování těchto událostí během anotace.}.

\begin{figure}[hbpt]
  \centering
  \includegraphics[width=0.9\textwidth]{./ch5-construction/img/energy_spec_normal.png}
  \caption{Průběh a spektrogram promluvy a vyznačenou energií promluvy.}
  \label{fig:construction:normal_speech}
\end{figure}

Dalším důvodem může být je nutnost správné artikulace. Při používání EL je to nezbytné, aby bylo produkované řeči alespoň trochu dobře rozumnět. A pokud se dobře artikuluje, tak není snadné mluvit rychle. Při nahrávání bylo také velmmi běžné, že v průběhu promluvy řečník udělal pauzu, aby mohl lépe umístit EL, protože jeho umístění má velký vliv na kvalitu produkované řeči. Nicméně je třeba říci, že tempo není a priory pro ASR systémy problém, protože různá délka fonémů je v relativně snadno modelována, např. v \textit{HMM} přechodem ze stavu do stejného stavu.

Dalším způsobem jak ukázat rozdíly mezi promluvou zdravého řečníka a řečníka s EL je srovnání ve frekvenční oblasti. Pro větší názornost jsou na obr. \ref{fig:construction:spectrogram} vedle sebe zobrazeno spektrum ukázková promluva zdravého řečníka (\ref{fig:construction:spectrogram:normal}) a toho s EL (\ref{fig:construction:spectrogram:el}). Obsah obou promluv je identický a přesto jsou obě spektra odlišná.

Prvním markantním rozdílem je mnohem větší zastoupení šumu v úsecích \uv{ticha} na obr. \ref{fig:construction:spectrogram:el}. To je nepochyně způsobnemo samotným EL, který řečník nevypíná mezi jednotlivými slovy. Na obr. \ref{fig:construction:el_speech} je to také zřetelně patrný, zejména na průběhu energie, šum zejména před prvním a druhým slovem promluvy. Zajímavá je přítomnost šumu v célém frekvenčním spektru, přestože EL produkuje konstantní buzení. Toto buzení je ve spektru, na obr. \ref{fig:construction:spectrogram:el}, viditelná jako výrazná souvislá linie v nízkých frekvencích. Přitomnost šumu ve vyšších frekvencích je způsobena umístěním mikrofonu, který je nalepen na pokožku a tím pádem snímá namodulované vibrace, přenášené měkkou tkání. Tento fakt se potvrdil v dalších etapách nahrávání (viz část \ref{chap:construction:normalization:corpus}), kde je použit studiový mikrofon vzdálený od úst minimálně 15 cm a tyto vibrace již nezaznamenává. Nicméně z pohledu použitelnosti nějakého budoucího systému je nezbytné počítat i se situací, kdy mikrofon bude zaznamenávat i vibrace přenášené tkání.

Dalším markantním rozdílem je absence vyšších frekvencí u většiny produkovaných fonémů. Vyjímku tvoří afrikáty $/c/$ a $/\check{c}/$, u kterých jsou hlasivky (u zdravého jedince) v klidu a vznikají uvolněním nahromaděného vzduchu v dutině ústní\footnote{Nahromadění vzduchu je realizováno přitisknutím jazyka k přední/zadní části horního patra.} \cite{Psutka2006}. V tomto případě není, u řečníka po TL, principiálně tento mechanizmus produkce těchto fonému ovlivněn. Problémem teoreticky může být zdroj vzduchu, jelikož jej z plic není možné dostat do dutiny ústní, ale jak už bylo zmíněno (a spektrogram to potvrzuje) zkušený uživatel EL se dokáže adaptovat.

Absence vyšších frekcencí se dá vysvětlit použitím EL, kde samotný EL má vždy konstantní frekvenci buzení a dále tím, že nedochází k modulaci v ve všech dutinách vokálního traktu. Nicméně nejdůležitější složky, zajišťující srozumitelnost, se vyskytují ve frekvenčním pásmu od 1 kHz do 3 kHz. Vyšší frekvence se a priory podílejí na zabarvení hlasu.

\begin{figure}[htpb]
  \centering
  \begin{subfigure}[b]{0.4\textwidth}
    \includegraphics[width=\textwidth]{./ch5-construction/img/spectrogram_normal.png}
    \caption{Zdravý řečník}
    \label{fig:construction:spectrogram:normal}
  \end{subfigure}
  %
  \begin{subfigure}[b]{0.4\textwidth}
    \includegraphics[width=\textwidth]{./ch5-construction/img/spectrogram_el.png}
    \caption{EL řečník}
    \label{fig:construction:spectrogram:el}
  \end{subfigure}
  \caption{Spektrogram promluvy \uv{Akcie Komerční banky} dvou řečníků.}
  \label{fig:construction:spectrogram}
\end{figure}

Dalším způsobem jak porovnat řeč zdravého řečníka a tím s EL je pomocí analýzy jednotlivých fonémů. Na obr. \ref{fig:construction:phonemes} jsou zobrazeny průběhy amplitudy v čase\footnote{Hodnoty času, na obr. \ref{fig:construction:phonemes}, odpovídají časům výskytu v původní promluvě.} pro fonémy $/k/$, $/m/$ a $/\check{c}/$. V případě $/k/$ a $/m/$ (1. a 2. průběh) se jedná o okluzivy, kde v prvním případě se jedná o neznělou plozivu a druhém o znělou plozivu. Tyto fonémy obecně vznikají uzavřením vydechovaného proudu vzduchu, pomocí artikulačních orgánů, což se projeví jako krátká pauza (tzv. okluze). Po té následuje náhlé jednorázové překážky a únik nahromaděného vzduchu, tzv. exploze \cite{Psutka2006}. Takto popsáno to samozřejmě funguje u zdravého jedince, ale u EL řečníka jde principiálně o stejný mechanizmus. S tím rozdílem, že vzduch nepochází z plic, ale z hltanu. Dalším rozdílem je samozřejmě absence hlasivek.

\begin{figure}[htpb]
  \centering
  \begin{subfigure}[b]{0.42\textwidth}
    \includegraphics[width=\textwidth]{./ch5-construction/img/phonemes_normal.png}
    \caption{Zdravý řečník}
    \label{fig:construction:phonemes:normal}
  \end{subfigure}
  %
  \begin{subfigure}[b]{0.42\textwidth}
    \includegraphics[width=\textwidth]{./ch5-construction/img/phonemes_el.png}
    \caption{EL řečník}
    \label{fig:construction:phonemes:el}
  \end{subfigure}
  \caption{Ukázky průběhů amplitudy pro fonémy $/k/$, $/m/$ a $/\check{c}/$.}
  \label{fig:construction:phonemes}
\end{figure}

Foném $/k/$ představuje zástupce neznělých fonémů, ty se vyznačují tím, že do jejich produkce nezasahují hlasivky, které jsou v klidu. Zdrojem buzení je tedy šum. Pokud se podíváme na průběh amplitudy v čase u zdravého řečníka (obr. \ref{fig:construction:phonemes:normal}), tak zde není vidět žádný periodický signál. Hlasivky jsou tedy opravdu v klidu. Oproti tomu u EL řečníka (obr. \ref{fig:construction:phonemes:el}) je jasně patrné, že je zde přítomno aktivní buzení vytvořené EL. Na obr. \ref{fig:construction:freq:k} je pak zobrazeno tzv. amplitudové spektrum, které znázorňuje vývoj amplitudy signálu ve frekcenci pro oba řečníky. V případě zdravého řečníka odpovídá vývoj očekávání tedy, že zde není žádná výrazná frekvence a také, že nedochází k výraznému útlumu. Přestože se v obou případech jedná o stejný foném, tak z časového i frekvešního průběhu amplitudy je zřejmé, že parametry signálu se u obou řečníku diametrálně liší.

\begin{figure}[hbpt]
  \centering
  \includegraphics[width=0.9\textwidth]{./ch5-construction/img/freq_analysis_(k).png}
  \caption{Vývoj amplitudy fonému $/k/$ ve frekvenci zdravého (horní) a EL (dolní) řečníka.}
  \label{fig:construction:freq:k}
\end{figure}

Jako druhý ukázkový foném slouží $/m/$. Opět se jedná o plozivu, ale v tomto případě o znělou. U těchto fonémů hrají velký vliv hlasivky, protože jsou zdrojem buzení. Z obr. \ref{fig:construction:phonemes:normal} je krásně zřetelné buzení ve formě perodického průběhu amplitudy. Narozdíl tomu, u EL řečníka (obr. \ref{fig:construction:phonemes:el}) je také vidět periodický signál, ale úplně jiného průběhu. Svým způsoběm dost podobný tomu, který je zřetelný u fonému $/k/$. Rozdíl je zřetelný i ve frekvenční oblasti (obr. \ref{fig:construction:freq:m}), kdy u EL řečnía nedochází útlumu ve střední oblasti frekvenčního spektra.

\begin{figure}[hbpt]
  \centering
  \includegraphics[width=0.9\textwidth]{./ch5-construction/img/freq_analysis_(m).png}
  \caption{Vývoj amplitudy fonému $/m/$ ve frekvenci zdravého (horní) a EL (dolní) řečníka.}
  \label{fig:construction:freq:m}
\end{figure}

Posledním úkázkovým fonémem je již zmiňované $/\check{c}/$. Jedná se o neznělý foném, který vzniká přiložením jazyku k zadní části horního patra. Tím je zadržen vzduch v dutině ústní a vzniká krátká pauza. Uvolněním pak dochází k explozi a vytvoření zvuku. Do produkce se nezapojují hlasivky a produkovaný zvuk by měl být dostatečně intenzivní, aby jej (v případě EL řečníka) tolik neovlivňoval EL. Tím pádem by měl být průběh signálu, u obou řečníků podobný, a to jak v časové, tak i ve frekvenční oblasti. Na obr. \ref{fig:construction:phonemes} a \ref{fig:construction:freq:c} je pak jasně vidět, že se jedná o platný předpoklad.

\begin{figure}[hbpt]
  \centering
  \includegraphics[width=0.9\textwidth]{./ch5-construction/img/freq_analysis_(c).png}
  \caption{Vývoj amplitudy fonému $/\check{c}/$ ve frekvenci zdravého (horní) a EL (dolní) řečníka.}
  \label{fig:construction:freq:c}
\end{figure}

Z doposud provedené analýzy plyne, že EL řeč je v mnoha charakteristikách odlišná od té produkované zdravým řečníkem. Zejména u porovnání ve frekvenční oblasti (obr. \ref{fig:construction:freq:k} a \ref{fig:construction:freq:m}) je to nejvíce patrné. Tento fakt nepochyně přispívá k tomu, že standardní obecné modely rozpoznávání řeči nedosahují takové přesnosti jako v případě bežné promluvy.

% \csvautotabular{./ch5-construction/test.csv}

% !TEX root = ../thesis.tex
\section{Akustický model}
\label{chap:construction:am}

TBD

% !TEX root = ../thesis.tex
\section{Jazykový model}
\label{chap:construction:lm}

% !TEX root = ../thesis.tex
\section{Aplikace obecného systému rozpoznávání a dosažené výsledky}
\label{chap:construction:results}

Z provedené analýzy plyne, že získaný EL korpus je odlišný od \uv{standardního} řečového korpusu, který se běžně používá  k~trénování obecných akustických modelů.
Tyto modely jsou nezávislé na řečníkovi a vyznačují se robustností.
Je tedy otázka, zda bude obecný akustický model schopen pracovat s~EL daty.

Za účelem ověření hypotézy byl vytvořen TDNN akustický model (popsaný v~části \ref{chap:asr:acoustic:DNN}), který byl natrénován daty z korpusu čítajícího $1000$ hodin promluv od 3951 řečníků.
Celkový počet HMM stavů je \textbf{3599}.

Jazykový model je založen na trigramech slov a  k~jeho natrénování posloužil textový korpus čítající velké množství novinových článků, webových reportáží, filmových titulků a dalších textových záznamů.
Slovník jazykového modelu čítá více než $1$ milion unikátních slov.

Testovacím vstupem vytvořeného ASR systému jsou data z EL korpusu.
Celková slovní přesnost, počítaná podle vzorce (\ref{eq:asr:decoding:acc}), dosáhla hodnoty $18,49~\%$\footnote{Dosaženo na state-of-the-art ASR systému v~době psaní práce. V~době vytvoření EL korpusu (kolem roku 2011) převládaly HMM-GMM akustické modely. Tento systém dosáhl přesnosti na slovech $12,59~\%$.}.
Dosažený výsledek zřetelně ilustruje odlišnost EL domény, protože obecný na řečníkovi nezávislý systém není schopen obstojně rozpoznat EL promluvu.

Pokud jsou  k~natrénování akustického modelu (taktéž využívajícího TDNN síť) použita pouze data\footnote{Korpus byl náhodně rozdělen na trénovací a testovací sadu v~poměru $90~\%\ (10h41m44s)$ trénovací a $10~\%\ (1h00m58s)$ testovací sada. Toto rozděleni je použito ve všech experimentech.} z EL korpusu, tak výsledná slovní přesnost dosáhla hodnoty $83,33~\%$, opět počítáno podle vzorce (\ref{eq:asr:decoding:acc}).
Jazykový model je identický jako v~případě obecného systému.
Dosažený výsledek demonstruje výhodu vytvoření individuálního modelu z EL nahrávek.
Zároveň ukazuje schopnost akustického modelu přizpůsobit se specifikům EL řeči.
Přestože je výstup individuálního modelu výrazně přesnější než obecného modelu zpracovávající EL promluvy, tak zdaleka nedosahuje výsledků srovnatelných s~nejlepšími ASR systémy, které jsou schopny v~ideálních podmínkách dosahovat na srovnatelné úloze více než $90~\%$ slovní přesnosti.


%Z dosaženého výsledku je patrné, že EL doména je diametrálně odlišná od běžné řeči, pro které jsou ASR systémy vytvářeny. Navíc, pokud se vezme v~potaz náročnost získání potřebných dat pro natrénování obecného modelu, tak se jako jediná schůdná varianta jeví vytváření individuálních modelů pro každého řečníka. To znamená, že model je trénovaný pouze z dat odpovídající konkrétnímu řečníkovi a často i účelu použití. K vytvoření takového modelu je zapotřebí řádově méně dat, při dosažení podobného výkonu. Stinnou stránkou je případná menší robustnost modelu. Čistě logicky tento model bude fungovat pouze s~konkrétním řečníkem a ještě jen v~situacích, které odpovídají trénovacím datům. U řečníků s~EL může navíc hrát velký vliv samotný EL. Již při nahrávání se ukázalo, že jeho pozice může nepříznivě ovlivnit kvalitu řeči. Tento problém by však neměl významně ovlivňovat kvalitu modelu, protože tento fenomém je obsažen v~datech.

\subsection{Hledání optimálních parametrů baseline modelu}
\label{chap:construction:results:baseline}

%Co se však ukázalo jako potencionálně problematické, je stabilita parametrů produkované řeči v~dlouhodobém časovém úseku. Více o tomto problému pak v~části \ref{chap:construction:normalization:quality}. K zodpovězení nejdůležitější otázky, jestli takový model vůbec může fungovat, stačí získaná data z první etapy nahrávání a ta obsahují řeč s~relativně konzistentními parametry.

V rámci ověřování funkčnosti individuálního modelu je vhodné otestovat různé kombinace parametrů modelu za účelem nalezení optimálních hodnot.
Hlavními uvažovanými hyperparametry jsou vzorkovací frekvence audio nahrávek a počet HMM stavů.
Pořízené nahrávky mají vzorkovací frekvenci rovnu $44,1\ kHz$.
Pro úlohu rozpoznávání EL řeči je to však zbytečně vysoká frekvence, protože nejhodnotnější informace je u EL řeči obsažena ve frekvenčním pásmu do $4\ kHz$.
Vyšší frekvence ovlivňuje především zabarvení hlasu a další individuální charakteristiky \cite{Psutka2006}.
Samotná EL řeč má svá specifika, proto se otevírá otázka volby vhodné vzorkovací frekvence.
% však obsahuje věci, které běžná řeč neobsahuje a tak je otázka, zde vhodnější vzorkovací frekvence rovna $8\ kHz$ nebo lépe $16\ kHz$.

Jak bylo řečeno výše, počet stavů modelu ovlivňuje množství modelovaných trifónů.
Čím více akustických jednotek je modelováno, tím více musí mít HMM model unikátních stavů.
Od počtu unikátních stavů se odvíjí rozsah množiny trénovacích dat potřebný pro  natrénování robustního modelu.
Pro účely výběru optimální struktury modelu bylo uvažováno $1024,\ 2048$ a $4096$ stavů.
%Jen pro vysvětlení je dobré zmínit, že HMM stav představuje model jedné uvažované akustické jednotky (nebo skupiny jednotek s~podobnými parametry). Počet stavů nám tedy říká, kolik takových jednotek model dokáže rozlišit. Čím více stavů, tím více jednotek (menších skupin) je modelováno. Teoreticky tak model s~více stavy je lepší. Nicméně  k~natrénování jednotky je potřeba určité množství dat, a tím pádem je pro model s~vyšším počtem stavů logicky potřeba větší množství trénovacích dat. Samozřejmě fonetická sada neobsahuje $4096$ fonémů, neobsahuje ani $1024$ fonémů\footnote{Ve skutečnosti obsahuje 42 českých fonémů.}. U těchto modelů se pak používá nějaký druh \textit{n-gramové} reprezentace fonémů, nejčastěji pak trifóny.
Jestliže byly uvažovány dvě hodnoty vzorkovací frekvence $8\ kHz$ a $16\ kHz$, bylo natrénováno $6$ modelů.
K vytvoření akustických modelů byl použit nástroj HTK-Toolkitu ve verzi v3.4., který je určen  k~vytváření HMM modelů.
Při trénování byl nejprve vytvořen monofónový akustický model s~jedním Gaussiánem pro každý stav.
Ten slouží jako základ pro trifónové modely.
Výsledný trifónový model využívá směs Gaussiánů v~souladu s~částí \ref{chap:asr:acoustic:GMM}.
Celý proces trénování je znázorněn na obr. \ref{fig:construction:results:baseline:hmm:training}.

\begin{figure}[hbpt]
  \centering
  \includegraphics[width=0.9\textwidth]{./ch5-construction/img/hmm-training.pdf}
  \caption{Princip trénování HMM-GMM modelu.}
  \label{fig:construction:results:baseline:hmm:training}
\end{figure}

Cílem popsané procedury je nalezení vhodných parametrů baseline modelu, zejména pak toho akustického.
K tomu je potřeba minimalizovat vliv jazykového modelu.
Z tohoto důvodu je použit speciální zerogramový jazykový model.
Standarně se předpokládá, že základní jednotkou jazykového modelu je slovo.
Pro ně jsou počítány četnosti z trénovacích dat a vytvořen model.
Obecně ale není nutné, aby základní jednotkou bylo slovo. V~případě experimentů s~EL řečí je mnohem výhodnější vytvořit model, jehož základní jednotkou je foném, protože právě foném je výstupem akustického modelu.
Výstupem ASR systému tak bude sekvence fonémů, která z~praktického pohledu není úplně užitečná, ale pro testování vlastností akustického modelu se hodí dokonale.
V rámci této práce je takovýto model označován jako fonémový zerogramový jazykový model.
Velikost slovníku tohoto modelu odpovídá velikosti fonémové sady, a tedy lze pravděpodobnost výskytu libovolného fonému stanovit jako $P(w_n) = 1/N$\footnote{Označení $w_n$ může evokovat použití slov u jazykového modelu. Změna písmena by však mohla vést ke zmatení čtenáře, protože by byla použita nestandardní notace. Z tohoto důvodu je i pro speciální model použito označení fonémů $w_n$.}, kde $N=40$.
Výsledná přesnost je tedy závislá pouze na akustickém modelu.
Akustická data byla parametrizována pomocí MFCC s~26 filtry, 12 kepstrálními koeficienty a energií.
Příznakový vektor obsahuje první i druhou derivaci těchto koeficientů.
Podrobněji je parametrizace popsána v~části \ref{chap:asr:parametrization:hearing}.


% Funkce ASR systému lze popsat rovnící

% \begin{equation}
%   \argmax_W p\left(W | O\right) = \argmax_W p\left(O | W\right) p\left(W\right),
% \end{equation}

%\noindent kde $\boldsymbol{O}$ reprezentuje sekvenci akustických příznaků a $W$ výstupní sekvenci znaků\footnote{Znakem tu může být myšleno písmeno, případně slovo.}. $P\left(O | W\right)$ je pravděpodobnost generování korektní pozorované sekvence, tedy korektní  k~akustickému modelu ASR systému. Pravděpodobnost $P\left(W\right)$ je apriorní pravděpodobnost konkrétní sekvence znaků $W$, jinými slovy jazykový model. K získání výsledků je tedy potřeba mít i tento model.

%Ten však není níjak ovlivněn řečníkem (pouze doménou použití systému) a není jej třeba upravovat pro potřeby řečníka s~EL. Cílem experimentu je ověření funkčnosti ASR a nalezení optimálních parametrů akustického modelu. Z tohoto důvodu je potřeba co nejvíce eliminovat vliv jazykového modelu na celkovém výkonu ASR systému. Jak bylo zmíněno, funkcí $p\left(W\right)$ je určení nejpravděpodobnější sekvence znaků. Pravděpodobnostní rozložení je získáno z velkého množství trénovacích textů. Toto natrénované rozložení by však velmi ovlivnilo výsledky experimentů, a proto je použít zerogramový monofónový model. Ten se vyznačuje tím, že všechny prvky slovníku mají stejnou pravděpodobnost rovnu $P(w_n) = \frac{1}{N}$, kde $N$ je počet položek ve slovníku. Monofónový model je navíc zvolen z toho důvodu, že fonetická sada je známa a obsahuje malý počet jednotek. Z pohledu jazykového modelu má libovolný výstup z akustického modelu stejnou pravděpodobnost výskytu, a tím pádem se jazykový model nijak nepřispívá  k~celkové kvalitě ASR systému.

V~tab. \ref{tab:construction:experiment:gmm} jsou uvedeny dosažené výsledky HMM-GMM modelů.
Opět se potvrdilo, že individuální ASR systém s~EL daty může fungovat.
Pokud dosažené výsledky porovnáme s~výsledky obecného modelu ($Acc_{word} = 18,49~\%)$, tak i zde je vidět rapidní nárůst přesnosti ($78,63~\%$ u nejhoršího individuálního GMM modelu).
Ze získaných výsledků je zřejmé, že volba vzorkovací frekvence má významný dopad na přesnost rozpoznávání.
Při použití vyšší z uvařovaných vzorkovacích frekvencí, zde $16\ kHz$, je vhodnější, protože oproti $8\ kHz$ bylo dosaženo zlepšení přesnosti o~$1,41~\%$ absolutně, tedy téměř $7~\%$ relativně.
Dodatečné experimenty ukázaly, že použití vyšší vzorkovací frekvence než $16\ kHz$ přinese jen zanedbatelné zlepšení.

Dále se ukázalo, že počet stavů nehraje tak zásadní roli při posuzování kvality akustického modelu, jako vzorkovací frekvence.
Z testované množiny počtu stavů dosáhl nejlepšího výsledku model, který měl maximálně $4096$ stavů, nicméně oproti modelu s~$1024$ stavy je nárůst přesnosti pouze $0,4~\%$ absolutně v~případě vzorkovací frekvence $16\ kHz$, což není tak významné.
Logicky se nabízí otázka, proč nezkusit ještě více stavů.
Odpověď se skrývá ve skutečném počtu stavů modelu s~maximálním počtem $4096$ stavů.
Slovíčko \uv{maximálním} je zde podstatné.
Algoritmus trénování akustického modelu se snaží rozdistribuovat všechny možné akustické jednotky (v~tomto případě trifóny) do maximálního počtu stavů.
Pokud chceme dosáhnout menšího počtu stavů než jednotek, dochází  k~určité formě shlukování (za pomocí fonetického rozhodovacího stromu) \cite{Holmes2001}.
Pokud je  k~dispozici dostatek dat  k~natrénování konkrétního shluku, je tento shluk akceptován, pokud jich není dostatečné množství, je tento shluk spojen s~jiným, který je svými parametry nejblíže.
V případě, že je  k~dispozici dostatek dat k~natrénování maximálního počtu stavů má model tento počet stavů.
Pokud není dostatek dat, může mít model méně stavů.
U modelu s~maximálním počtem $4096$ stavů je skutečný počet stavů $3257$, tzn. i kdyby se trénoval model s~$8192$ stavy, tak by se tato hodnota změnila jen velmi málo.

%Pro doplnění, fonémový akustický model dosáhl přesnosti $54,49~\%$ pro $8\ kHz$ a $62,30~\%$ pro $16\ kHz$.

\begin{table}[htpb]
  \centering
  \def\arraystretch{1.5}
  \pgfplotstabletypeset[
    col sep=semicolon,
    string type,
    columns/model/.style={column name={Počet HMM stavů}, column type={c}},
    columns/8k/.style={column name={8 kHz}, column type={r}},
    columns/16k/.style={column name={16 kHz}, column type={r}},
    every head row/.style={
      before row={
        \toprule & \multicolumn{2}{c}{$Acc_{p}\ [\%]$} \\
      },
      after row={
        \cmidrule(r){1-1}
        \cmidrule(l){2-3}
      }
    },
    every last row/.style={after row={\bottomrule}},
  ]{./ch5-construction/tabs/01-frequency.csv}
  \caption{Vliv frekvence na kvalitu modelu.}
  \label{tab:construction:experiment:gmm}
\end{table}

Hledání optimálních parametrů baseline modelu bylo realizováno na přelomu let $2013$ a $2014$.
V tuto dobu byly stále dominantní GMM modely.
Z tohoto důvodu byl později tento experiment zopakován s~HMM-DNN akustickým modelem.
Vstupem modelu byla stejná MFCC parametrizace s~26 filtry, 12 kepstrálními koeficienty plus energie, delta a delta-delta příznaky.
Tato parametrizace je provedena na mikrosegmentu $t$ a jeho okolí $t-5$ a $t+5$.
Každý mikrosegment má délku $10\ ms$.
Samotná síť se skládá z 6 vrstev, každá s~4096 neurony, výstupní vrstva je typu softmax s~dimenzí rovnou počtu HMM stavů.
Dosažené výsledky jsou v~tab. \ref{tab:construction:experiment:dnn}.
Z nich je patrné, že nalezené optimální hyperparametry jsou shodné i při použití DNN.
Nicméně je zde i zřejmý důvod následné dominance HMM-DNN modelů.
Fungují totiž výrazně lépe.
Pouhou náhradou GMM za DNN bylo dosaženo zlepšení.
Hodnoty přesnosti baseline modelu jsou tedy pro GMM $Acc_{p}^{GMM} = 81,20~\%$ a pro DNN $Acc_{p}^{DNN} = 85,23~\%$. To odpovídá zlepšení o~$4~\%$ absolutně oproti nejlepšímu GMM výsledku, viz tab. \ref{tab:construction:experiment:gmm} a \ref{tab:construction:experiment:dnn}.

%experiment byl realizován na přelomu let $2013$ a $2014$, kdy ještě $ASR$ modelům dominovaly \textit{HMM-GMM} modely, byl později zopakován s~\textit{HMM-DNN} modely, které dosahují ještě vyšších přesností.

%Více o \textit{HMM-DNN} v~části \ref{chap:construction:normalization:corpus}. Výsledky těchto modelů jsou v~tab. \ref{tab:construction:experiment:dnn}, z nich je vidět, že i v~této oblasti neuronové sítě jasně dominují.

\begin{table}[htpb]
  \centering
  \def\arraystretch{1.5}
  \pgfplotstabletypeset[
    col sep=semicolon,
    string type,
    columns/model/.style={column name={Počet HMM stavů}, column type={c}},
    columns/8k/.style={column name={8 kHz}, column type={r}},
    columns/16k/.style={column name={16 kHz}, column type={r}},
    every head row/.style={
      before row={
        \toprule & \multicolumn{2}{c}{$Acc_{p}\ [\%]$} \\
      },
      after row={
        \cmidrule(r){1-1}
        \cmidrule(l){2-3}
      }
    },
    every last row/.style={after row={\bottomrule}},
  ]{./ch5-construction/tabs/01-frequency_dnn.csv}
  \caption{Vliv frekvence na kvalitu modelu využívajícího DNN.}
  \label{tab:construction:experiment:dnn}
\end{table}

\subsection{Redukce fonetické sady}
\label{chap:construction:results:reduction}

Při mluvení je elektrolarynx permanentně zapnutý, a to i v~případě neznělých fonémů.
Toto má vliv na jejich průběh, podrobněji v~části v~\ref{chap:construction:analysis}.
Nabízí se tak předpoklad, že všechny neznělé fonémy mají podobu znělých párových fonémů, a tím pádem je možné redukovat fonetickou sadu.
V důsledku redukce fonetické sady by došlo ke snížení perplexity modelu.
Rozhodnutí, zda se jedná o variantu slova obsahující znělý nebo neznělý foném, je pak přenecháno jazykovému modelu.

K ověření tohoto předpokladu je potřeba provést experiment.
Myšlenka experimentu je jednoduchá.
Je potřeba natrénovat několik modelů lišících se pouze tím, jaký fonetický pár (viz tab. \ref{tab:construction:reduction:pairs}) je použit pro redukci fonetické sady.
V rámci experimentu jsou uvažovány tyto případy:

\begin{itemize}
  \item \textit{Baseline} - standardní model s~plnou fonetickou sadou.
  \item $/f/ \rightarrow /v/$ - foném $/f/$ je nahrazen fonémem $/v/$.
  \item $/k/ \rightarrow /g/$ - foném $/k/$ je nahrazen fonémem $/g/$.
  \item $/s/+/\check{s}/ \rightarrow /z/+/\check{z}/$ - foném $/s/$ $\left(/\check{s}/\right)$ je nahrazen fonémem $/z/$ $\left(/\check{z}/\right)$.
  \item $/t/+/\text{\textit{ť}}/ \rightarrow /d/+/\text{\textit{ď}}/$ - foném $/t/$ $\left(/\text{\textit{ť}}/\right)$ je nahrazen fonémem $/d/$ $\left(/\text{\textit{ď}}/\right)$.
  \item \textit{Náhrada všech} - všechny uvažované neznělé fonémy jsou nahrazeny znělým ekvivalentem.
\end{itemize}

\begin{table}[htpb]
  \centering
  \def\arraystretch{1.5}
  \pgfplotstabletypeset[
    col sep=comma,
    string type,
    columns/unvoiced/.style={column name={Neznělé fonémy}, column type={c}},
    columns/voiced/.style={column name={Znělé fonémy}, column type={c}},
    every head row/.style={
      before row={\toprule},
      after row={\cmidrule(r){1-1}\cmidrule(l){2-2}}
    },
    every last row/.style={after row={\bottomrule}},
  ]{./ch5-construction/tabs/phonemes_pairs.csv}
  \caption{Korespondující páry fonémů.}
  \label{tab:construction:reduction:pairs}
\end{table}

\noindent Pro porovnání jsou stejné modely vytvořeny i pro zdravého řečníka.
U něj by při libovolné redukci fonetické sady mělo dojít ke zhoršení oproti \textit{baseline} modelu.

K natrénování akustických modelů byly použity korpusy čítající 5000 vět\footnote{Pro oba řečníky byly použity stejné věty pocházející z databáze popsané v~\cite{Radova2000}.}, což představuje více než 10 hodin řeči jak pro zdravého, tak i EL řečníka.
Akustická data byla parametrizována pomocí MFCC s~26 filtry a 12 kepstrálními koeficienty a energií.
Dále vektor parametrů obsahuje delta a delta-delta příznaky.
To dohromady tvoří vektor 39 příznaků pro každých 10 ms nahrávky \cite{Psutka2007}.

V rámci experimentu byly otestovány dva přístupy.
%vzájemně se lišící řečovou jednotkou.
V prvním případě se jednalo o monofónový akustický model, v~druhém případě byl testován trifónový.
U obou přístupů je řečová jednotka reprezentována pětistavovým HMM-GMM modelem se spojitou výstupní pravděpodobnostní funkcí pro každý stav.
% , viz \ref{chap:asr:acoustic:HMM}.
Pro určení optimálních parametrů modelu pro EL byly využity poznatky uvedené v~ části \ref{chap:construction:results:baseline}.
Pro zdravého řečníka je pro každou část experimentu vytvořeno několik modelů lišících se počtem stavů a Gaussovských směsí.
Všechny akustické modely jsou natrénovány pomocí HTK-Toolkitu v3.4.
Celkem bylo vytvořeno 24 akustických modelů, 12 pro EL řečníka (6 monofónových a 6 trifónových) a 12 pro zdravého řečníka.

Pro otestování modelů byla vytvořena testovací sada čítající 500 vět náhodně vybraných z původních korpusů (pro oba řečníky stejná).
Testovací sada tak představuje přibližně 1 hodinu řeči pro každého řečníka.
V rámci tohoto experimentu jsou uvažovány dva jazykové modely

\begin{enumerate}
  \item \textit{zerogramový jazykový model} - v~tomto případě mají všechna slova v~modelu stejnou pravděpodobnost výskytu $P_r(w_n|w_1,\dots,w_{n-1}) = \frac{1}{N}$, kde $N$ je počet slov ve slovníku. Konkrétně pro $N = 2885$, odpovídá perplexita modelu hodnotě $2885$. Testovací slovník je vytvořen z testovací sady, model tedy neobsahuje OOV\footnote{Out-of-vocabulary (OOV) - slova, která nejsou obsažena ve slovníku jazykového modelu.}.
  \item \textit{trigramový jazykový model} - u tohoto modelu odpovídá pravděpodobnost výskytu následujícího slova $P_r(w_n|w_1,\dots,w_{n-1})~=~P(w_n|w_{n-2}, w_{n-1})$. K získání $P(w_n|w_{n-2}, w_{n-1})$ posloužil SRILM Toolkit s~Kneser-Ney vyhlazováním\footnote{Vyhlazování slouží  k~vyřešení problému s~OOV, kdy trénovací data neobsahovala tato OOV, a proto není  k~dispozici $p(w_n|w_{n-2}, w_{n-1})$.} \cite{Stolcke2002}, které se podle \cite{Prazak2008} ukázalo jako optimální pro tyto typy modelů. Jako trénovací data byly použity texty z novinových článků, webových stránek a přepisů televizních pořadů. Celkem model obsahuje 360 tisíc nejvíce frekventovaných slov. OOV bylo $3,8 \%$ a perplexita $3380$.
\end{enumerate}

\noindent v~kombinaci s~vytvořenými akustickými modely to představuje 4 dílčí experimenty.
Jen pro doplnění je nutné poznamenat, že přesnost modelů je vyhodnocována na slovech.
%, oproti fonémovému baseline modelu v~\ref{chap:construction:results:baseline}.

V tab. \ref{tab:construction:reduction:01} a \ref{tab:construction:reduction:02} jsou shrnuty dosažené výsledky pro monofónový akustický model a zerogramový jazykový model, resp. trigramový jazykový model.
V obou případech se potvrdilo očekávané chování přesnosti modelu u zdravého řečníka.
Redukcí fonetické sady je omezena komplexita modelu, a tím pádem dochází ke zhoršení přesnosti.
Překvapivé mohou být horší dosažené výsledky jak u zdravého, tak i EL řečníka, viz tab. \ref{tab:construction:reduction:02}.
Toto chování může být vysvětleno vyšší perplexitou trigramového jazykového modelu v~kombinaci s~relativně jednoduchým monofónovým akustickým modelem.

U EL řečníka je vidět dílčí zlepšení u dvou modelů (tab. \ref{tab:construction:reduction:01}), resp. jednoho modelu v~případě trigramového modelu (tab. \ref{tab:construction:reduction:02}).
Ve většině případů však redukce fonetické sady vedla ke zhoršení přesnosti.
Při použití trigramového jazykového modelu došlo obecně ke zlepšení výkonu systému i přesto, že tento model má vyšší perplexitu.
% To nasvědčuje tomu, že pro tuto úlohu je monofónový akustický model příliš jednoduchý.

\begin{table}[htpb]
  \centering
  \def\arraystretch{1.5}
  \pgfplotstabletypeset[
    col sep=semicolon,
    string type,
    columns/model/.style={column name=Model, column type={c}},
    columns/normal/.style={column name={Zdravý}, column type={r}},
    columns/el/.style={column name={EL}, column type={r}},
    every head row/.style={
      before row={
        \toprule & \multicolumn{2}{c}{$Acc_{w}\ [\%]$} \\
      },
      after row={
        \cmidrule(r){1-1}
        \cmidrule(l){2-3}
      }
    },
    every last row/.style={after row={\bottomrule}},
  ]{./ch5-construction/tabs/reduction_01.csv}
  \caption[Vliv redukce na přesnost ASR s~mono. AM. a zerogram LM.]{Vliv redukce fonetické sady na přesnost ASR systému s~monofoním akustickým a zerogramovým jazykovým modelem ($N=2885$) pro zdravého a EL řečníka.}
  \label{tab:construction:reduction:01}
\end{table}

\begin{table}[htpb]
  \centering
  \def\arraystretch{1.5}
  \pgfplotstabletypeset[
    col sep=semicolon,
    string type,
    columns/model/.style={column name=Model, column type={c}},
    columns/normal/.style={column name={Zdravý}, column type={r}},
    columns/el/.style={column name={EL}, column type={r}},
    every head row/.style={
      before row={
        \toprule & \multicolumn{2}{c}{$Acc_{w}\ [\%]$} \\
      },
      after row={
        \cmidrule(r){1-1}
        \cmidrule(l){2-3}
      }
    },
    every last row/.style={after row={\bottomrule}},
  ]{./ch5-construction/tabs/reduction_02.csv}
  \caption[Vliv redukce na přesnost ASR s~mono. AM. a 3-gram LM (N=360k).]{Vliv redukce fonetické sady na přesnost ASR systému s~monofonovým akustickým a trigramovým jazykovým modelem obsahujícím 360 tisíc slov pro zdravého a EL řečníka.}
  \label{tab:construction:reduction:02}
\end{table}


V tab. \ref{tab:construction:reduction:03} a \ref{tab:construction:reduction:04} jsou pak uvedeny výsledky pro trifónový akustický model se zerogramovým resp. trigramovým jazykovým modelem.
Stejně jako u předchozích dvou experimentů, tak i zde, je vidět, že redukce fonetické sady vede u zdravého řečníka vždy ke zhoršení přesnosti modelu.
Z uvedených výsledků je patrné, že trifónový akustický model dosahuje výrazně lepších výsledků než monofonový model.
Lze se domnívat, že zhoršení kvality rozpoznávání u EL řečníka v~tab. \ref{tab:construction:reduction:03} je s~nejvyšší pravděpodobností způsobeno využitím fonetických stromů v~AM, protože v~trénovací sadě není dostatek dat pro všechny možné varianty trifónů.
Tím pádem model pro určité trifóny vrací špatné sekvence znaků.
Zerogramový jazykový model pak nedokáže pomoci, protože všechna slova mají stejnou pravděpodobnost výskytu $P_r(w_n|w_1,\dots,w_{n-1}) = \frac{1}{2885}$.
Tím pádem dochází  k~rozpoznání špatného slova, tedy ke snížení celkové přesnosti rozpoznávání.
Tuto domněnku potvrzuje rapidní zlepšení v~případě trigramového jazykového modelu (tab. \ref{tab:construction:reduction:04}), kde již jazykový model významně přispívá  k~přesnosti modelu.

\begin{table}[htpb]
  \centering
  \def\arraystretch{1.5}
  \pgfplotstabletypeset[
    col sep=semicolon,
    string type,
    columns/model/.style={column name=Model, column type={c}},
    columns/normal/.style={column name={Zdravý}, column type={r}},
    columns/el/.style={column name={EL}, column type={r}},
    every head row/.style={
      before row={
        \toprule & \multicolumn{2}{c}{$Acc_{w}\ [\%]$} \\
      },
      after row={
        \cmidrule(r){1-1}
        \cmidrule(l){2-3}
      }
    },
    every last row/.style={after row={\bottomrule}},
  ]{./ch5-construction/tabs/reduction_03.csv}
  \caption[Vliv redukce na přesnost ASR s~trif. AM. a zerogram. LM.]{Vliv redukce fonetické sady na přesnost ASR systému s~trifónovým akustickým a zerogramovým jazykovým modelem pro zdravého a EL řečníka.}
  \label{tab:construction:reduction:03}
\end{table}

\begin{table}[htpb]
  \centering
  \def\arraystretch{1.5}
  \pgfplotstabletypeset[
    col sep=semicolon,
    string type,
    columns/model/.style={column name=Model, column type={c}},
    columns/normal/.style={column name={Zdravý}, column type={r}},
    columns/el/.style={column name={EL}, column type={r}},
    every head row/.style={
      before row={
        \toprule & \multicolumn{2}{c}{$Acc_{w}\ [\%]$} \\
      },
      after row={
        \cmidrule(r){1-1}
        \cmidrule(l){2-3}
      }
    },
    every last row/.style={after row={\bottomrule}},
  ]{./ch5-construction/tabs/reduction_04.csv}
  \caption[Vliv redukce na přesnost ASR s~trif. AM. a 3-gram LM (N=360k).]{Vliv redukce fonetické sady na přesnost ASR systému s~trifónovým akustickým a trigramovým jazykovým modelem s~360 tisíc slov pro zdravého a EL řečníka.}
  \label{tab:construction:reduction:04}
\end{table}

U obou experimentů s~trifónovým akustickým modelem došlo ke zlepšení u dvou modelů (tab. \ref{tab:construction:reduction:03} a \ref{tab:construction:reduction:04}). Naproti tomu redukce fonetické sady vedla ve většině případů stejně jako v~případě monofónového modelu ke zhoršení přesnosti rozpoznávání.


Ze získaných výsledků je možné usoudit, že redukce fonetické sady může vést ke zlepšení přesnosti.
Nicméně předpoklad, že všechny neznělé fonémy jsou shodné se svými znělými ekvivalenty se ukázala jako mylná.
Zároveň také není možné říci, že náhrada např. dvojice $/s/$ a $/\check{s}/$  za každých okolností povede  k~lepším výsledkům.
Při hlubší analýze se ukázalo, že velmi záleží na kontextu daného fónemu, ten totiž velmi ovlivňuje jeho podobu.
Řeč představuje spojitou formu signálu a při vyslovování různých slov obsahujících stejný foném s~odlišným okolím  může dojít  k~odchylkám například v~artikulaci, příkladem může být dvojice slov \textit{hrad} a \textit{had}.
Toto pozorování ověřil i dodatečný experiment, ve kterém nebyla provedena náhrada $/s/$ za $/z/$ pouze pro případ trifónu \textit{b-s+t}, který je například ve slově \textit{obstát}.
Díky vynechání tohoto jediného trifónu byla výsledná nejlepší přesnost u trifónového akustického modelu $83,39~\%$ (původně $83,28~\%$) v~případě zerogramového jazykového modelu a $88,44~\%$ (původně $88,31~\%$) v~případě trigramového modelu.
Přestože se jedná o marginální zlepšení, tak ho bylo docíleno jedním trifónem.
Bohužel určení toho, jaké trifóny vynechat z nahrazování není v~žádném případě triviální úloha.

Zajímavý je také rozdíl mezi přesností modelu pro zdravého řečníka a EL řečníka.
Přestože se v~obou případech jedná o individuální modely šité \uv{na míru} řečníkovi, tak průměrný rozdíl v~přesnosti rozpoznávání činí $6,24~\%$ absolutně, resp. $40,38~\%$ relativně.
To značí, že je potřeba se zabývat myšlenkou úpravy akustického modelu za účelem dosažení lepších výsledků.
V~ideálním případě dokonce srovnatelných s~modely pro zdravé řečníky.

Naopak očekávaným výsledkem bylo zhoršení přesnosti rozpoznávání u zdravého řečníka ve všech případech redukce fonetické sady.
Dále se potvrdilo, že komplexnější trifónový model dosahuje ve většině případů lepších výsledků.
To je nepochybně způsobeno faktem, že každý foném je v~případě trifónového modelu modelován pomocí více HMM stavů, než je tomu u monofónového modelu.

% %!TEX root = ../thesis.tex
\section{Rehabilitace hlasu po totální laryngektomii}
\label{chap:cause:treatment}

Nesporná výhoda totální laryngektomie neoddiskutovatelně spočívá v~odstranění
primárního nádorového onemocnění. Následky operace však způsobí obrovský
zásah do kvality života pacienta. Okem nejviditelnější změnu představuje
přítomnost otvoru na krku po provedeném operačním zákroku a s~ním spojený způsob dýchání.
Do uměle vyvedené průdušnice je přímo vdechován vzduch z okolního prostředí a nedochází tak  k~jeho přirozeně probíhající filtraci, ohřevu ani přirozenému zvlhčování. Toto má za následek vyšší náchylnost pacientů  k~respiračním onemocněním.

Lze se domnívat, že pro samotného pacienta je jedním z nejobtížnějších úkolů vypořádat se s~trvalou
ztrátou vlastního hlasu. Z toho důvodu se již samotný autor operačního zákroku doktor
Billroth zaobíral otázkou rehabilitace hlasu. Jeho první pokusy s~kovovou
tracheostomickou kanylou sice umožňovaly pacientovi hovořit, ale svou
konstrukcí pacienta spíše ohrožovaly na životě, proto se častěji využívala
metoda tzv. jícnového hlasu \cite{Sebova-Sedenkova2006}. Ve stejnou
dobu, tedy začátkem minulého století, se začaly objevovat také %první interní a externí
hlasové aparáty. V~současnosti se rehabilitace hlasu provádí s~využitím: %následujících metod:

\begin{itemize}
  \item \textbf{foniatrických metod}, mezi které patří metoda jícnového hlasu a využití elektrolarynxu,
  \item \textbf{chirurgicko-protetickým způsobem}, který spočívá v~opětovném propojení průdušnice a jícnu,
  \item \textbf{vytvoření hrtanu podobných struktur chirurgickým způsobem},
  \item \textbf{transplantace hrtanu}.
\end{itemize}

\noindent Může se zdát, že je  k~dispozici relativně široká škála možností, jak pacientovi vrátit schopnost vyjadřování pomocí mluvené řeči.
Ovšem je nutné si uvědomit, že výběr konkrétní metody závisí na stavu
a možnostech pacienta.
Jinými slovy, ne každá metoda je vhodná pro každého pacienta a žádná z~metod není univerzální.
Souhrnný přehled spolu s~hlavními výhodami a nevýhodami jednotlivých rehabilitačních přístupů je uveden v~tab. \ref{tab:treatment:summary}.

% \subsection{Foniatrické metody} % (fold)
% \label{chap:cause:treatment:foniatric}

% Nevyhnutelným důsledkem odstranění hrtanu je ztráta hlasu. Neznamená to ale, že by byla
% úplně eliminována schopnost produkovat řeč. V~procesu vytváření hlasu zastává
% odstraněný orgán pouze (i když velmi zásadní) roli generátoru zvuku. Zbylé
% orgány, jako například hrdelní, nosní a ústní dutina, zůstávají nedotčeny a mohou i
% nadále plnit svou funkci. Logicky se tak nabízí myšlenka nahradit chybějící
% zdroj zvuku jiným. Tento princip využívají metoda jícnového hlasu a produkce řeči s~využitím
% elektrolarynxu.

% \subsubsection{Jícnový hlas} % (fold)
% \label{chap:cause:treatment:foniatric:esophageal}

% První zmínky o využívání jícnového hlasu se datují do roku 1922, kdy prof. MUDr. Miloslav Seeman
% \cite{seeman1922speech} potvrdil domněnku, že funkci štěrbiny mezi hlasivkami (rima
% glottidis) přebírá tzv. pseudoglottis, která se vytváří na úrovni horního
% jícnového svěrače, a vypracoval metodiku vytváření jícnového
% hlasu.

% Princip tvorby jícnového hlasu spočívá v~tom, že se vzduch neplní do plic, ale do jícnu.
% Takto si pacient připravuje potřebný vzduch  k~následné
% eruktaci\footnote{eruktace - latinsky název pro proces říhání (popřípadě
% krkání), při kterém dochází  k~úniku plynů pocházejících ze žaludku dutinou
% ústní.} vzduchu a produkci řeči. Vlastní jícnový hlas vzniká na přechodu
% jícnu a hypofaryngu (spodní část hltanu). Následně v~oblasti horního jícnového
% zúžení dochází  k~rozkmitání sliznice a podslizniční vrstvy, a tedy  k~produkci zvuku,
% který je následně modulován stejně jako v~případě přirozené produkce řeči.
% Princip tvorby \uv{základního} tónu jícnového hlasu je znázorněn na obr.
% \ref{fig:cause:treatment:esophageal}.

% \begin{figure}[htb]
%   \begin{center}
%     \includegraphics[width=0.6\linewidth]{ch3-cause/figures/esophageal}
%     \caption[Princip produkce jícnového hlasu.]{Princip produkce jícnového hlasu. Průchodem vzduchu přes zúžení vzniká základní tón jícnového hlasu.}
%     \label{fig:cause:treatment:esophageal}
%   \end{center}
% \end{figure}

% Na základě toho, jakým způsobem je jícen plněn vzduchem, rozlišujeme metodu \textbf{aspirační} a metodu \textbf{injekční}. Zatímco aspirační metoda spočívá v~plnění jícnu vzduchem pomocí polykání, u injekční metody je za tímto účelem využíván kořen jazyka, pomocí něhož je vzduch vtlačován do jícnu. Následný princip produkce hlasu je pro oba případy shodný. Injekční princip plnění jícnu vzduchem je využíván pacienty, kterým byla
% při laryngektomii odstraněna jazylka. V~těchto případech nelze jícen naplnit aspirační cestou.

% Proces učení jícnového hlasu by měl začít co možná nejdříve po operaci. Pokud
% je to možné, tak se s~výukou začíná ještě za pobytu pacienta na ORL klinice
% nebo krátce po propuštění. V~první fázi se pacient učí pouze slabiky
% sestávající z explosivy a souhlásky. Postupně se však přidávají slabičné
% shluky, které sice nedávají smysl, ale pomáhají v~osvojení potřebné techniky.
% V případě úspěšného zvládnutí se přistupuje  k~nácviku frází a souvislé řeči.
% Potřebnou dobu  k~nácviku jícnového hlasu nelze přesně určit, protože je
% závislá na mnoha faktorech. V~literatuře se uvádí, že pro úspěšné osvojení techniky jícnového hlasu je potřeba 30 až 50 hodin velmi intenzivního tréninku \cite{Brown2003}. Míra úspěšnosti nácviku srozumitelného hlasu se uvádí v~rozsahu 14\%-75\%. %Takto obrovský rozsah značí o mnoha faktorech, které mohou ovlivnit úspěšné osvojení jícnového hlasu.
% Mezi možné příčiny, které mohou významně ovlivnit zvládnutí techniky jícnového hlasu, patří
% fyziologické potíže, anatomické problémy, psychologické problémy, nebo jednoduše
% neadekvátní podpora při řečové terapii \cite{Brown2003}. Významnou roli také
% hraje snaha a odhodlání samotného pacienta.

% Nepopiratelnou výhodou využití metody jícnového hlasu je, že při rehabilitaci hlasu není pacient závislý na lékaři. Navíc po chirurgickém zásahu je zajištěno permanentní oddělení dýchacích a
% polykacích cest, je tedy eliminováno riziko vniknutí potravy do dýchacích cest.
% Další nespornou výhodou pro pacienty, kteří ovládají techniku jícnového hlasu, je bezesporu to, že mají při mluvení
% volné obě ruce. Za nevýhodu lze obecně považovat nižší srozumitelnost produkovaného hlasu.
% Ta je způsobena tzv. \uv{břišním} zabarvením, ke kterému zcela jistě při produkci řeči prostřednictvím metody jícnového hlasu  dochází, a dále nízkou intenzitou hlasu a krátkou výdrží mluvčího při tvorbě tónu. Za negativum se dá také považovat množství pacientem vynaloženého úsilí potřebného  k~osvojení
% techniky. Velmi často se také mluvčí ostýchají jícnový hlas používat, protože
% mají pocit, že je společensky nevhodné dorozumívat se formou blízkou říhání. Z
% tohoto důvodu se odhaduje, že v~běžném životě využívá jícnový hlas pouze 20\% -30\%
% pacientů, kteří se začali tuto techniku učit \cite{Hradecka2007}.

% subsubsection jícnový_hlas (end)

% \subsubsection{Využití elektrolarynxu} % (fold)
% \label{chap:cause:treatment:foniatric:elektrolarynx}

% Rehabilitace hlasu za pomoci přídavného zařízení, nazývaného elektrolarynx, se řadí mezi tzv. elektromechanické
% metody. Princip metody spočívá v~přikládání speciálního zařízení obsahujícího generátor zvuku do oblasti spodiny úst. Generátor zvuku, který je součástí elektrolarynxu, napomáhá přenosu zvuku a vibrací nejen do dutiny ústní, ale i do přilehlých artikulačních orgánů. Následně je pacient schopen běžným způsobem artikulovat, a tedy i hovořit. Princip vytváření hlasu pomocí elektrolarynxu je ukázán na obr. \ref{fig:cause:treatment:electrolarynx}.

% \begin{figure}[htb]
%   \begin{center}
%     \includegraphics[width=0.6\linewidth]{ch3-cause/figures/electrolarynx}
%     \caption[Princip rehabilitace hlasu pomocí elektrolarynxu.]{Princip rehabilitace hlasu pomocí elektrolarynxu.}
%     \label{fig:cause:treatment:electrolarynx}
%   \end{center}
% \end{figure}

% Takto generovaná řeč se vyznačuje několika charakteristickými rysy. V~první
% řadě řeč budí velmi mechanický dojem. Důvodem je samozřejmě samotný
% elektrolarynx, který lze označit za elektromechanický generátor zvuku s
% konstantním buzením. Proto je také základní frekvence produkovaného hlasu víceméně konstantní a řečník tak má velmi omezené možnosti jak řeč emotivně zabarvovat. V~průběhu času se sice objevily snahy o ovlinění základní frekvence produkované řeči \cite{Kikuchi2004, Uemi1994, Goldstein2004}, kterou by bylo možné vyvolat prostřednictvím změny budící frekvence elektrolarynxu, ale ty ztroskotaly, jelikož jako velmi obtížné se jeví nalezení vhodného mechanismu, pomocí kterého by bylo docíleno optimální změny fundamentální frekvence mluvené řeči s~ohledem na to, co chce řečník říci. Dalším charakteristickým rysem hlasu produkovaného pomocí elektrolarynxu je jeho nižší srozumitelnost. K tomuto jevu dochází v~důsledku přítomnosti zvukového podkresu generovaného přístrojem. Srozumitelnost se navíc snižuje s~rostoucím okolním hlukem, proto se velmi často stává, že posluchač, který se s~takto produkovanou řečí setkává poprvé, není schopen promluvám plně porozumět.

% Naproti tomu, za výhody, kterými tato metoda rehabilitace hlasu disponuje, lze označit možnost rychlého
% osvojení schopnosti opětovně produkovat řeč. Navíc je vhodná pro téměř všechny pacienty
% postižené ztrátou hlasu v~důsledku odstranění hrtanu popř. jeho poškození.
% Za nevýhodu lze obecně pokládat kvalitu produkované řeči, tedy monotonní a
% mechanicky znějící hlas. Za jistý způsob omezení lze považovat nutnost držení elektrolarynxu, jakožto přídavného zařízení, při mluvení.

% \subsection{Chirurgicko-protetická metoda} % (fold)
% \label{chap:cause:treatment:tracheo}

% Další možností, kterou lze využít pro rehabilitaci hlasu, je po odstranění hrtanu chirurgicky vytvořit průchod
% mezi průdušnicí a jícnem tak, aby u tracheostomovaného člověka mohl opětovně proudit vzduch z plic do úst.
% Princip metody spočívá v~tom, že pacient při výdechu zneprůchodní operativně vytvořený otvor (stoma) v~oblasti krku. Vzduch tak projde skrz fistuli do oblasti jícnu, naráží do jeho stěn a dochází  k~jeho rozvibrování. Následně jsou tyto vibrace modulovány pomocí artikulačních orgánů a vzniká řeč.

% Historicky první zmínka o vytvoření fistule\footnote{fistule (česky píštěl) je abnormální
% otvor mezi dvěma dutými orgány, nebo mezi dutým orgánem a kůží.} mezi
% průdušnicí a jícnem pochází z~roku 1932. V~tomto roce doktor Guttman poprvé
% vytvořil tracheoezofageání shunt\footnote{shunt - kanál, kterým je tekutina
% odkloněna z přirozené dráhy. Tento kanál může být vytvořen chirurgicky nebo pomocí syntetické trubice. }.
% Hlavní snahou chirurgů bylo vytvoření bezpečné, správně nasměrované píštěle
% umožňující tvorbu hlasu. Bohužel v~mnoha případech byly zákroky doprovázené
% vážnými komplikacemi (infekcemi, zápaly či těžkým krvácením).
% Mimo to se museli operatéři vypořádat se zajištěním stálosti
% vytvořeného otvoru tak, aby jím neprotékaly tekutiny špatným směrem a
% nedocházelo  k~jejich zatékání do dýchacích cest. Jelikož se jednalo o velmi
% náročné operační postupy a bylo s~nimi spojeno velké množství rizik, došlo v
% 80.~letech 20.~století  k~jejich ústupu. Svou renesanci zažily s~návrhem první protézy v~podobě jednocestného ventilu, který zajišťoval pouze jednosměrný
% průchod tekutin skrze píštěl, jak je ilustrováno na obr.
% \ref{fig:cause:treatment:shunt}. První komerčně dostupná protéza se objevila
% v~80.~letech 20.~století v~USA.

% \begin{figure}[htb]
%   \begin{center}
%     \includegraphics[width=0.6\linewidth]{ch3-cause/figures/te-shunt}
%     \caption[Průchod vzduchu tracheoezofageální protézou.]{Průchod vzduchu tracheoezofageální protézou.}
%     \label{fig:cause:treatment:shunt}
%   \end{center}
% \end{figure}

% Na používané protézy jsou kladené
% přísné nároky a musí vyhovovat určitým požadavkům. Musí být vyrobeny
% z~biokompatibilního materiálu, který odolává biodegradaci. Tím je zaručena jejich
% dlouhodobá trvanlivost a správná funkce. V~neposlední řadě by měla být protéza
% samofixační a snadno vyměnitelná.
%  K zajistění jejich správné funkce je zapotřebí navrhnout je tak, aby
% tlak potřebný  k~otevření faryngoezofageálního segmentu byl co nejnižší. To by mělo zajišťovat pacientovi vytvářet
% plynulou řeč. První vyráběné protézy se ale vyznačovaly tím, že tlak potřebný pro jejich otevření byl příliš vysoký a omezovaly tak množinu potencionálních pacientů. Nejmodernější protézy se již vyznačují
% velmi nízkým otevíracím fonačním tlakem.

% V praxi se používá několik druhů protéz. Hlavním rozdílem mezi nimi však je,
% zda se pacient přímo účastní výměny ventilu, jehož fundamentální funkcí je
% vytvoření průchodu pro vzduch proudící z průdušnice do jícnu. U protéz, které
% jsou vyměňovány operačně, se doba používání pohybuje od 3 do 6 měsíců. Tento
% interval velmi významně ovlivňuje tvorba biofilmu na povrchu náhrady. K tvorbě
% dochází následkem přímého kontaktu protézy s~tělními tekutinami a potravou.
% Rychlost tvorby biofilmu ovlivňuje tvar a materiál, ze kterého je náhrada
% vytvořena \cite{Leunisse2001}. U typů, které si nositel může měnit sám, se
% předpokládá, že budou čištěny nebo měněny přibližně jednou za dva týdny. Na obr. \ref{fig:cause:treatment:prosthesis} jsou
% ukázány některé typy užívaných protéz.

% \begin{figure}[htb]
%   \begin{center}
%     \includegraphics[width=0.9\linewidth]{ch3-cause/figures/te-protezy}
%     \caption[Ilustrace používaných TE protéz.]{Ilustrace používaných TE protéz (a) Gronigenova nízkotlaká protéza, (b) Provox2 a (c) Blom-Singer protéza.}
%     \label{fig:cause:treatment:prosthesis}
%   \end{center}
% \end{figure}

% Samotný zákrok zavedení protézy je možné provést současně s~výkonem totální
% laryngektomie (tzv. primární zavedení hlasové protézy) nebo až po zotavení
% pacienta z~náročné léčby nádorového onemocnění (tzv. sekundární zavedení).
% Primární zavedení umožňuje začít s~hlasovou rehabilitací krátce po odstranění
% hrtanu. Zároveň pacient nemusí v~krátké době podstupovat druhou operaci, při
% které by se vkládal jednocestný ventil do vytvořené fistule.

% V praxi se ukázalo, že úspěšnost rehabilitace pomocí této metody je více než 80~\%
% \cite{Slavicek2002}. Důležitým faktorem, stejně jako u metody jícnového hlasu, je
% funkčnost faryngoezofageálního segmentu. Významnou roli hraje také otvírací tlak horního
% jícnového svěrače. Hlas tvořený protézou se vyznačuje vysokou kvalitou, dobrou
% srozumitelností, individuálním zabarvením a relativně dlouhou fonační dobou
% dosahující průměrně 20 sekund \cite{Saito2003}. Oproti jícnovému hlasu není potřeba
% tak intenzivní edukace pacienta k~plnému osvojení hlasu. V~současnosti se
% jedná o nejpoužívanější metodu rehabilitace hlasu.

% \subsection{Hrtanu podobné struktury} % (fold)
% \label{chap:cause:tratment:structure}

% S rozvojem mikrovaskulárních\footnote{mikrovaskulární - část oběhového systému
% složeného z nejmenších cév, jako jsou kapiláry, žilky aj.} transplantátů se
% začaly objevovat postupy, které umožňovaly rehabilitovat hlas pouze pomocí
% chirurgického zákroku. Tyto techniky jsou založeny na myšlence permanentního spojení hypofaryngu
% s tracheou pomocí vlastní tkáně pacienta.

% První metodu tohoto druhu představil v~roce 1984 doktor Ehrenberger
% \cite{Kramp2009}, který popsal tzv. \uv{\textbf{řečový sifón}}
% (angl. \textbf{speech siphon}). Jedná se o dvakrát esovitě zahnuté spojení mezi hrtanem a hltanem
% vytvořené z části tenkého střeva zvané lačník (jejunum). Dvojesovité zahnutí je provedeno tak, aby
% bylo minimalizováno riziko sekundární aspirace. Schéma \uv{řečového sifónu}
% podle Ehrenberga je znázorněno na obr. \ref{fig:cause:treatment:microvascular}
% A. Již na první pohled je zřejmé, že se jedná o velmi náročný chirurgický
% zákrok. První články publikované autorským kolektivem prezentovaly velmi dobré
% funkční výsledky metody. Podle \cite {Sebova-Sedenkova2006} bylo touto metodou doposud
% operováno přibližně 60 pacientů.

% V roce 1990 byla popsána laryngoplastika podle Hagena. V~tomto případě je za účelem vytvoření ústrojí zvaného \textbf{neolarynx} používán štěp z předloktí. Neoglottis je vyztužen chrupavkou
% a překrývá vchod do neolaryngu tak, aby nedocházelo  k~sekundární aspiraci. Vnitřek neolaryngu je kryt kůží.
% Laryngoplastika podle Hagena je znázorněna na obr.
% \ref{fig:cause:treatment:microvascular} B. Doposud bylo touto metodou operováno přibližně 300
% pacientů \cite{Sebova-Sedenkova2006}.

% \begin{figure}[htb]
%   \begin{center}
%     \includegraphics[width=0.9\linewidth]{ch3-cause/figures/microvascular}
%     \caption[Schéma \uv{řečového sifónu} a laryngoplastiky.]{A) Schéma \uv{řečového sifónu} tak jak jej představil Ehrenberg. B) Laryngoplastika podle Hagena.}
%     \label{fig:cause:treatment:microvascular}
%   \end{center}
% \end{figure}

% Bohužel v~současné době tyto metody nenacházejí širší uplatnění. Hlavním důvodem je
% náročnost realizace operačních postupů, kvůli které se velmi
% těžko prosazují na dalších pracovištích. Dalším aspektem, který limituje tyto
% metody, je vliv na samotného pacienta. Metody předpokládají provedení dalšího chirurgického
% zákroku, který představuje pro pacienta další zátěž a mohou ho provázet nemalé komplikace.
% I přes nedostatky těchto metod je
% pochopitelná snaha lékařů o intenzivní výzkum v~této oblasti. Při úspěšné
% léčbě je pacient schopen produkovat hlas velmi dobré kvality a ve většině
% případů nepotřebuje žádnou péči ze strany lékařů.

% \subsection{Transplantace hrtanu}
% \label{chap:cause:treatment:transplantation}

% Nejkomplexnější možnost rehabilitace hlasu představuje transplantace hrtanu.
% Pokud je úspěšná, přebírá transplantovaný orgán plně funkci původního
% orgánu a velmi významně zvyšuje šance pacienta na plné zotavení bez trvalých
% následků. Jedná se ale o velmi náročný chirurgický zákrok, protože je potřeba provést reinervaci a obnovení cévního zásobení implantátu.

% Přestože první zmínky o možnosti provedení transplantace hrtanu se objevují již v~60.~letech 20.~století\footnote{Vůbec první úspěšná transplantace
% orgánu (ledvin) se uskutečnila v~roce 1954.}, byla první transplantace tohoto druhu provedená
% až profesorem Marshallem Stromem v~roce 1998 \cite{Narula2011},
% Prvním pacientem, který podstoupil zmiňovaný chirurgický zákrok, byl čtyřicetiletý muž z USA.
% K~laryngektomii v~jeho případě vedla motocyklová nehoda, při které došlo  k~rozdrcení hrtanu. Před transplantací pacient využíval 20 let  k~produkci řeči elektrolarynx. Dárcem orgánu byl taktéž čtyřicetiletý muž, který zemřel na následky prasknutí mozkového aneurysmatu. Příjemce transplantátu již třetí den po operaci promluvil (vyslovil
% anglické slovo \uv{hello}). Přibližně po 36 měsících od transplantace byl
% produkovaný hlas kvalitativně srovnatelný s~hlasem zdravého člověka. U této transplantace se však nepodařilo dosáhnout kompletní reinervace, proto nebylo možné zajistit bezproblémové dýchání a bylo proto nutné
% ponechat tracheostomii. I přes tuto skutečnost měl zákrok významný podíl na zvýšení kvality jeho života
% \cite{Strome2001}. Doposud poslední úspěšně vykonaná transplantace hrtanu byla
% dle dostupných zdrojů provedena v~říjnu 2010.

% Jedním z hlavních důvodů tak nízkého počtu úspěšně provedených zákroků je, že se jedná o transplantaci dárcovského orgánu. Pacientům jsou podávány medikamenty zabraňující odmítnutí dárcovského orgánu (imunosupresiva), která ale v~současné době nelze podávat pacientům trpícím rakovinným onemocněním, protože významně zvyšují riziko opětovného rozšíření rakoviny \cite{Narula2011}. Ve výjimečných případech lze tuto metodu zvažovat u pacientů, kteří trpěli benigními nádory a minimálně 5~let u nich nedošlo  k~recidivě.
% Poslední výzkumy však naznačují, že by v~dohledné době v~této oblasti mohlo dojít  k~pokroku, v~důsledku kterého by bylo možno provést transplantaci hrtanu i u lidí, kteří trpěli rakovinným onemocněním.

% \subsection{Shrnutí} % (fold) \label{sub:treatment:summary}

% Rehabilitaci hlasu u pacientů, kteří prodělali chirurgické odstranění hrtanu, je ve
% vyspělých zemích věnována značná pozornost, jelikož následky této operace velmi významně ovlivňují kvalitu jejich života. Léčení jedinci se musí vyrovnat primárně se ztrátou hlasu. V~některých případech dochází po operaci i ke ztrátě čichu a pacienti jsou náchylnější ke vzniku respiračních onemocnění.
% Tato situace je již sama o sobě velmi náročnou psychickou zkouškou. Neméně významnou
% roli sehrává i fyzická odlišnost a z~toho pramenící psychická zátěž, které je pacient vystaven
% po absolvované léčbě.

% V~současnosti je za účelem rehabilitace hlasu využíváno několik metod, jejichž
% Souhrnný přehled spolu s~hlavními výhodami a nevýhodami jednotlivých rehabilitačních přístupů je uveden v~tab. \ref{tab:treatment:summary}.
% U většiny pacientů je hlas rehabilitován pomocí metody tracheoezofageálního píštěle,
% která je principiálně založena na metodě jícnového hlasu. O úspěchu rehabilitace, stejně jako u jícnového hlasu, tak
% především rozhodují vlastnosti faryngoezofageálního segmentu. Pokud si pacient
% není schopen osvojit techniku jícnového hlasu, případně nemá voperován píštěl, je
% pro rehabilitaci hlasu použit elektrolarynx. Za nejkomplexnější ze zmíněných postupů se dá považovat úplná transplantace hrtanu, která
% řeší víceméně všechny problémy spojené s~odstraněním hrtanu. Bohužel tento
% zákrok je velmi náročný a vhodný pouze pro omezený okruh pacientů.

\newcolumntype{b}{X}
\newcolumntype{s}{>{\hsize=.5\hsize}X}

\begin{table}[ht]
  \centering
  \begin{tabularx}{1.0\textwidth}{L{1.2} L{0.6} L{1.1} L{1.1}}
    & \textbf{Kvalita} & \textbf{Výhody} & \textbf{Nevýhody} \\
    \toprule \\ [-1.75ex]

    \textbf{Tracheoezofageální píštěl} & Vysoká & Vysoká míra osvojení, dlouhá fonační doba & Zanášení píštěle a s~ním spojené čištění, případně dodatečná lékařská péče \\
    \midrule \\ [-1.75ex]

    \textbf{Jícnový hlas} & Dobrá & Volné ruce při mluvení, není potřeba dodatečné lékařské péče & Velmi náročná metoda  k~naučení, nepřirozený hlas \\
    \midrule \\ [-1.75ex]

    \textbf{Elektrolarynx} & Nízká & Snadné  k~naučení & Monotonní až robotický hlas, nutné nosit externí elektrické zařízení \\
    \midrule \\ [-1.75ex]

    \textbf{Hrtanu podobné struktury} & Vysoká & Nezávislost pacienta na pravidelné lékařské péči & Velmi náročná chirurgická procedura, která pacienta vystavuje dalším možným rizikům  \\
    \midrule \\ [-1.75ex]

    \textbf{Transplantace hrtanu} & Velmi vysoká & Transplantovaný hrtan přejímá funkci odstraněného orgánu & Velmi náročná chirurgická procedura, která je vhodná jen pro malé procento pacientů \\
  \end{tabularx}

  \caption{Přehled dostupných metod rehabilitace hlasu. \label{tab:treatment:summary}}
\end{table}

% Přestože lékařská věda v~současnosti disponuje možnostmi jak rehabilitovat hlas a lékaři se snaží využívané postupy neustále zdokonalovat, zůstává zde otevřený prostor pro inovace, a tím zlepšení kvality života lidí postižených ztrátou hrtanu.

